\documentclass[12pt,a4paper]{ctexbook}
\usepackage{geometry}
\geometry{margin=2.2cm}
\usepackage{setspace}
\setstretch{1.35}
\usepackage{hyperref}
\title{全宋词选·ci.song.7000}
\author{古典文献整理}
\date{}
\begin{document}
\maketitle
\tableofcontents
\newpage
\section{好事近}
\textbf{蔡伸}\par\vspace{0.5em}
十幅健帆风，天意巧催行客。\par
极目五湖云浪，泛满空秋色。\par
玉人应怪误佳期，凝恨正脉脉。\par
锦鳞为传尺素，报兰舟消息。\par

\section{卜算子}
\textbf{蔡伸}\par\vspace{0.5em}
风雨送春归，寂寞花空委。\par
枝上红稀地上多，万点随流水。\par
翠黛敛春愁，照影临清泚。\par
应念韶华惜蕣颜，洒遍胭脂泪。\par

\section{卜算子}
\textbf{蔡伸}\par\vspace{0.5em}
小阁枕清流，一霎莲塘雨。\par
风递幽香入槛来，枕簟全无暑。\par
遐想似花人，阅岁音尘阻。\par
物是人非空断肠，梦入芳洲路。\par

\section{卜算子}
\textbf{蔡伸}\par\vspace{0.5em}
玉斧斫冰轮，中有乘鸾女。\par
鬓乱钗横襟袖凉，只恐轻飞举。\par
青冥缥缈间，自有吹箫侣。\par
不向巫山十二峰，朝暮为云雨。\par

\section{卜算子}
\textbf{蔡伸}\par\vspace{0.5em}
前度月圆时，月下相携手。\par
今夜天边月又圆，夜色如清昼。\par
风月浑依旧。\par
水馆空回首。\par
明夜归来试问伊，曾解思量否。\par

\section{卜算子}
\textbf{蔡伸}\par\vspace{0.5em}
重重雪外山，渺渺烟中路。\par
路转山横无尽愁，正是分携处。\par
望极锦中书，肠断鱼中素。\par
锦素沈沈两未期，鱼雁空相误。\par

\section{卜算子}
\textbf{蔡伸}\par\vspace{0.5em}
春事付莺花，曾是莺花主。\par
醉拍春衫金缕衣，只向花间住。\par
密意君听取。\par
莫逐风来去。\par
若是真心待于飞，云里千条路。\par

\section{小重山}
\textbf{蔡伸}\par\vspace{0.5em}
楼外江山展翠屏。\par
沈沈虹影畔，彩舟横。\par
一尊别酒为君倾。\par
留不住，风色太无情。\par
斜日半山明。\par
画栏重倚处，独销凝。\par
片帆回首在青冥。\par
人不见，千里暮云平。\par

\section{小重山}
\textbf{蔡伸}\par\vspace{0.5em}
澹澹秋容烟水寒。\par
楼高清夜永，倚阑干。\par
玉人不见坐长叹。\par
箫声远，明月满空山。\par
遐想绿云鬟。\par
青冥风露冷，独乘鸾。\par
别时容易见时难。\par
凭孤枕，聊复梦蝉娟。\par

\section{小重山}
\textbf{蔡伸}\par\vspace{0.5em}
流水桃花小洞天。\par
壶中春不老，胜尘寰。\par
霞衣鹤氅并桃冠。\par
新装好，风韵愈飘然。\par
功行满三千。\par
婴儿并姹女，炼成丹。\par
刘郎曾约共升仙。\par
十个月，养个小金坛。\par

\section{小重山}
\textbf{蔡伸}\par\vspace{0.5em}
楼上风高翠袖寒。\par
碧云笼淡日，照阑干。\par
绿杨芳草恨绵绵。\par
长亭路，何处认征鞍。\par
晓镜懒重看。\par
鬓云堆凤髻，任阑珊。\par
鸾衾鸳枕小屏山。\par
人如玉，忍负一春闲。\par

\section{踏莎行}
\textbf{蔡伸}\par\vspace{0.5em}
佩解江皋，魂消南浦。\par
人生惟有别离苦。\par
别时容易见时难，算来却是无情语。\par
百计留君，留君不住。\par
留君不住君须去。\par
望君频问梦中来，免教肠断巫山雨。\par

\section{踏莎行}
\textbf{蔡伸}\par\vspace{0.5em}
如是我闻，金仙出世。\par
一超直入如来地。\par
慈悲方便济群生，端严妙相谁能比。\par
四众归依，悉皆欢喜。\par
有情同赴龙华会。\par
无忧帐里结良缘，麽诃修哩修修哩。\par

\section{踏莎行}
\textbf{蔡伸}\par\vspace{0.5em}
客里光阴，伤离情味。\par
玉觞未举心先醉。\par
临岐莫怪苦留连，樯乌转处人千里。\par
恨写新声，云笺密寄。\par
短封难尽心中事。\par
凭君看取纸痕斑，分明总是离人泪。\par

\section{踏莎行}
\textbf{蔡伸}\par\vspace{0.5em}
落日归云，寒空断雁。\par
吴波浅淡山平远。\par
丹青写出在霜缣，佳人特地栽团扇。\par
渔艇孤烟，酒旗幽院。\par
些儿景趣君休羡。\par
五湖归去共扁舟，何如早早酬深愿。\par

\section{踏莎行}
\textbf{蔡伸}\par\vspace{0.5em}
水满青钱，烟滋翠葆。\par
残英满地无人扫。\par
先来羁思乱如云，无端更被春酲恼。\par
叠叠遥山，绵绵远道。\par
凭阑满目唯芳草。\par
莫惊青鬓点秋霜，卢郎已分愁中老。\par

\section{踏莎行}
\textbf{蔡伸}\par\vspace{0.5em}
玉质孤高，天姿明慧。\par
了无一点尘凡气。\par
白莲空殿锁幽芳，亭亭独占秋光里。\par
一切见闻，不可思议。\par
我今有分亲瞻礼。\par
愿垂方便济众生，他时同赴龙华会。\par

\section{定风波}
\textbf{蔡伸}\par\vspace{0.5em}
一曲骊歌酒一锺。\par
可怜分袂太匆匆。\par
百计留君留不住。\par
□去。\par
满川烟暝满帆风。\par
目断魂销人不见。\par
但见。\par
青山隐隐水浮空。\par
拟把一襟相忆泪。\par
试□。\par
云笺密洒付飞鸿。\par

\section{定风波}
\textbf{蔡伸}\par\vspace{0.5em}
老去情锺不自持。\par
篸花酌酒送春归。\par
玉貌冰姿人窈窕。\par
一笑。\par
清狂岂减少年时。\par
欲上香车俱脉脉，半帘花影月平西。\par
待得酒醒人已去。\par
凝伫。\par
断云残雨尽堪悲。\par

\section{点绛唇}
\textbf{蔡伸}\par\vspace{0.5em}
水绕孤城，乱山深锁横江路。\par
帆归别浦。\par
苒苒兰皋暮。\par
人在天涯，雁背南云去。\par
空凝伫。\par
凤楼何处。\par
烟霭迷津渡。\par

\section{点绛唇}
\textbf{蔡伸}\par\vspace{0.5em}
香雪飘零，暖风著柳笼丝雨。\par
恼人情绪。\par
春事还如许。\par
宝勒朱轮，共结寻芳侣。\par
东郊路。\par
乱红深处。\par
醉拍黄金缕。\par

\section{点绛唇}
\textbf{蔡伸}\par\vspace{0.5em}
背壁灯残，卧听檐雨难成寐。\par
井梧飘坠。\par
历历蛩声细。\par
数尽更筹。\par
滴尽罗巾泪。\par
如何睡。\par
甫能得睡。\par
梦到相思地。\par

\section{点绛唇}
\textbf{蔡伸}\par\vspace{0.5em}
月缺花残，世间乐事难双美。\par
夜来相对。\par
把酒弹轻泪。\par
一点情锺，销尽英雄气。\par
樊笼外。\par
五湖烟水。\par
好作扁舟计。\par

\section{点绛唇}
\textbf{蔡伸}\par\vspace{0.5em}
玉笋持杯，敛红颦翠歌金缕。\par
彩鸳戢羽。\par
未免群鸡妒。\par
我为情多，愁听多情语。\par
君休诉。\par
两心坚固。\par
云里千条路。\par

\section{点绛唇}
\textbf{蔡伸}\par\vspace{0.5em}
人面桃花，去年今日津亭见。\par
瑶琴锦荐。\par
一弄清商怨。\par
金日重来，不见如花面。\par
空肠断。\par
乱红千片。\par
流水天涯远。\par

\section{点绛唇}
\textbf{蔡伸}\par\vspace{0.5em}
梅雨初晴，画栏开遍忘忧草。\par
兰堂清窈。\par
高柳新蝉噪。\par
枕上芙蓉，如梦还惊觉。\par
匀妆了。\par
背人微笑。\par
风入玲珑罩。\par

\section{点绛唇}
\textbf{蔡伸}\par\vspace{0.5em}
帐外华灯，翠屏花影参差满。\par
锦衣香暖。\par
苦恨春宵短。\par
画角声中，云雨还轻散。\par
河桥畔。\par
月华如练。\par
回首成肠断。\par

\section{点绛唇}
\textbf{蔡伸}\par\vspace{0.5em}
绿萼冰花，数枝清影横疏牖。\par
玉肌清瘦。\par
夜久轻寒透。\par
忍使孤芳，攀折他人手。\par
人归后。\par
断肠回首。\par
只有香盈袖。\par

\section{点绛唇}
\textbf{蔡伸}\par\vspace{0.5em}
解绂朝天，满城桃李繁阴市。\par
彩舟难驻。\par
忍听骊歌举。\par
协赞中兴，圣意方倾注。\par
从今去。\par
五云深处。\par
稳步沙堤路。\par

\section{点绛唇}
\textbf{蔡伸}\par\vspace{0.5em}
云雨匆匆，洞房当日曾相遇。\par
暂来还去。\par
无计留春住。\par
宝瑟重调，静听莺弦语。\par
休轻负。\par
绮窗朱户。\par
好做风光主。\par

\section{昭君怨}
\textbf{蔡伸}\par\vspace{0.5em}
一曲云和松响。\par
多少离愁心上。\par
寂寞掩屏帷。\par
泪沾衣。\par
最是销魂处。\par
夜夜绮窗风雨。\par
风雨伴愁眠。\par
夜如年。\par

\section{醉落魄・一斛珠}
\textbf{蔡伸}\par\vspace{0.5em}
波纹如縠。\par
池塘雨后添新绿。\par
海棠初绽红生肉。\par
双燕归来，还认旧巢宿。\par
凝情凭暖阑干曲。\par
新愁无限伤心目。\par
谁人月下吹横玉。\par
惊起鸳鸯，飞去自相逐。\par

\section{醉落魄・一斛珠}
\textbf{蔡伸}\par\vspace{0.5em}
霜华摇落。\par
亭亭皓月侵朱箔。\par
梦回敧枕听残角。\par
一片寒声，风送入寥廓。\par
眼前风月都如昨。\par
独眠无奈情怀恶。\par
凭肩携手于飞约。\par
料想人人，终是赋情薄。\par

\section{醉落魄・一斛珠}
\textbf{蔡伸}\par\vspace{0.5em}
明眸秀色。\par
双蛾巧画春山碧。\par
盈盈标韵倾瑶席。\par
一见尊前，宛是旧相识。\par
深期密语虽端的。\par
良宵无奈成轻掷。\par
忍教只恁空相忆。\par
得入手来，无限好则剧。\par

\section{醉落魄・一斛珠}
\textbf{蔡伸}\par\vspace{0.5em}
阳关声咽。\par
清歌响断云屏隔。\par
溪山依旧连空碧。\par
昨日主人，今日是行客。\par
绿窗朱户应如昔。\par
回头往事成陈迹。\par
后期总便无端的。\par
月下风前，应也解相忆。\par

\section{极相思}
\textbf{蔡伸}\par\vspace{0.5em}
碧檐鸣玉玎珰。\par
金锁小兰房。\par
楼高夜永，飞霜满院，璧月沈缸。\par
云雨不成巫峡梦，望仙乡、烟水茫茫。\par
风前月底，登高念远，无限凄凉。\par

\section{极相思}
\textbf{蔡伸}\par\vspace{0.5em}
相思情味堪伤。\par
谁与话衷肠。\par
明朝见也，桃花人面，碧藓回廊。\par
别后相逢唯有梦，梦回时、展转思量。\par
不如早睡，今宵魂梦，先到伊行。\par

\section{玉楼春}
\textbf{蔡伸}\par\vspace{0.5em}
碧桃溪上蓝桥路。\par
寂寞朱门闲院宇。\par
粉墙疏竹弄清蟾，玉砌红蕉宜夜雨。\par
个中人是吹箫侣。\par
花底深盟曾共语。\par
人生乐在两知心，此意此生君记取。\par

\section{玉楼春}
\textbf{蔡伸}\par\vspace{0.5em}
星河风露经年别。\par
月照离亭花似雪。\par
宝钗鸾镜会重逢，花里同眠今夜月。\par
月华依旧当时节。\par
细把离肠和泪说。\par
人生只合镇长圆，休似月圆圆又缺。\par

\section{长相思}
\textbf{蔡伸}\par\vspace{0.5em}
我心坚。\par
你心坚。\par
各自心坚石也穿。\par
谁言相见难。\par
小窗前。\par
月婵娟。\par
玉困花柔并枕眠。\par
今宵人月圆。\par

\section{长相思}
\textbf{蔡伸}\par\vspace{0.5em}
锦衾香。\par
玉枕双。\par
昨夜深深小洞房。\par
回头已断肠。\par
背兰缸。\par
梦仙乡。\par
风撼梧桐雨洒窗。\par
今宵好夜长。\par

\section{长相思}
\textbf{蔡伸}\par\vspace{0.5em}
村姑儿。\par
红袖衣。\par
初发黄梅插稻时。\par
双双女伴随。\par
长歌诗。\par
短歌诗。\par
歌里真情恨别离。\par
休言伊不知。\par

\section{西地锦}
\textbf{蔡伸}\par\vspace{0.5em}
寂寞悲秋怀抱。\par
掩重门悄悄。\par
清风皓月，朱阑画阁，双鸳池沼。\par
不忍今宵重到。\par
惹离愁多少。\par
蓬山路杳，蓝桥信阻，黄花空老。\par

\section{归田乐}
\textbf{蔡伸}\par\vspace{0.5em}
风生苹末莲香细。\par
新浴晚凉天气。\par
犹自倚朱阑，波面双双彩鸳戏。\par
鸾钗委坠云堆髻。\par
谁会此时情意。\par
冰簟玉琴横，还是月明人千里。\par

\section{七娘子}
\textbf{蔡伸}\par\vspace{0.5em}
天涯触目伤离绪。\par
登临况值秋光暮。\par
手拈黄花，凭谁分付。\par
雍雍雁落蒹葭浦。\par
凭高目断桃溪路。\par
屏山楼外青无数。\par
绿水红桥，锁窗朱户。\par
如今总是销魂处。\par

\section{感皇恩}
\textbf{蔡伸}\par\vspace{0.5em}
酒晕衬横波，玉肌香透。\par
轻袅腰肢妒垂柳。\par
臂宽金钏，且是不干春瘦。\par
拈金双合字，无心绣。\par
鬒云半堕，金钗欲溜。\par
罗袂残香忍重嗅。\par
渡江桃叶，肠断为谁招手。\par
倚阑凝望久，眉空斗。\par

\section{感皇恩}
\textbf{蔡伸}\par\vspace{0.5em}
膏雨晓来晴，海棠红透。\par
碧草池塘袅金柳。\par
王孙何在，不念玉容消瘦。\par
日长深院静，帘垂绣。\par
璨枕堕钗，粉痕轻溜。\par
玉鼎龙涎记同嗅。\par
钿筝重理，心事谩凭纤手。\par
素弦弹不尽，眉峰斗。\par

\section{减字木兰花}
\textbf{蔡伸}\par\vspace{0.5em}
彤庭龙尾。\par
礼备天颜知有喜。\par
九奏初传。\par
耳冷人间十七年。\par
盈成持守。\par
仁德如春渐九有。\par
三辅名州。\par
好整笙歌结胜游。\par

\section{减字木兰花}
\textbf{蔡伸}\par\vspace{0.5em}
船回沙尾。\par
几误红窗听鹊喜。\par
尺素空传。\par
转首相逢又隔年。\par
寒灯独守。\par
玉笋持杯宁复有。\par
秀水南州。\par
徒使幽人作梦游。\par

\section{减字木兰花}
\textbf{蔡伸}\par\vspace{0.5em}
多情多病。\par
玉貌疲来愁览镜。\par
门掩东风。\par
零落桃花满地红。\par
重帘不卷。\par
愁睹杏梁双语燕。\par
强拂瑶琴。\par
一曲幽兰泪满襟。\par

\section{减字木兰花}
\textbf{蔡伸}\par\vspace{0.5em}
金风玉露。\par
喜鹊桥成牛女渡。\par
天宇沈沈。\par
一夕佳期两意深。\par
琼签报曙。\par
忍使飙轮容易去。\par
明日如今。\par
想见君心似我心。\par

\section{减字木兰花}
\textbf{蔡伸}\par\vspace{0.5em}
锦屏人醉。\par
玉暖香融春有味。\par
今日兰舟。\par
魂梦还随绿水流。\par
高城望断。\par
无奈城中人不见。\par
斜倚妆楼。\par
恨入眉峰两点愁。\par

\section{渔家傲}
\textbf{蔡伸}\par\vspace{0.5em}
烟锁池塘秋欲暮。\par
细细前香，直到双栖处。\par
并枕东窗听夜雨。\par
偎金缕。\par
云深不见来时路。\par
晓色朦胧人去住。\par
香覆重帘，密密闻私语。\par
日断征帆归别浦。\par
空凝伫。\par
苔痕绿印金莲步。\par

\section{西楼子・相见欢}
\textbf{蔡伸}\par\vspace{0.5em}
楼前流水悠悠。\par
驻行舟。\par
满目寒云衰草、使人愁。\par
多少恨，多少泪，谩迟留。\par
何似蓦然拚舍、去来休。\par

\section{西楼子・相见欢}
\textbf{蔡伸}\par\vspace{0.5em}
红靴玉带葳蕤。\par
翠绡衣。\par
并辔垂鞭妆影、照清溪。\par
长亭路。\par
停骑处。\par
晚凉时。\par
空有许多明月、伴双栖。\par

\section{御街行}
\textbf{蔡伸}\par\vspace{0.5em}
东君不锁寻芳路。\par
曾是莺花主。\par
有情风月可怜宵，犹记绿窗朱户。\par
十年空想，春风面杳，无计凭鳞羽。\par
凄凉怀抱今如许。\par
天与重相遇。\par
不应还向楚峰前，朝暮为云为雨。\par
算来各把，平生分付，也不是、恶著处。\par

\section{上阳春・蓦山溪}
\textbf{蔡伸}\par\vspace{0.5em}
好在章台杨柳。\par
不禁春瘦。\par
淡烟微雨曲尘丝，锁一点、眉头皱。\par
忆自灞陵别后。\par
青青依旧。\par
万丝千缕太多情，忍攀折、行人手。\par

\section{临江仙}
\textbf{蔡伸}\par\vspace{0.5em}
繁杏枝头蜂蝶乱，香风阖坐微闻。\par
靓妆浓艳任东君。\par
无情风雨，春事已平分。\par
珍重主人留客意，夜阑秉烛开尊。\par
何须歌韵遏行云。\par
羽觞交劝，挥麈细论文。\par

\section{临江仙}
\textbf{蔡伸}\par\vspace{0.5em}
昨夜中秋今夕望，十分桂影团圆。\par
玉人相对绿尊前。\par
素娥有恨，应是妒蝉娟。\par
人静小庭风露冷，歌声特地清圆。\par
醉红醺脸髻鬟偏。\par
翠裙轻皱，端的为留仙。\par

\section{临江仙}
\textbf{蔡伸}\par\vspace{0.5em}
帘幕深深清昼永，玉人不耐春寒。\par
镂牙棋子缕金圆。\par
象盘雅戏，相对小窗前。\par
隔打直行尖曲路，教人费尽机关。\par
局中胜负定谁偏。\par
饶伊使幸，毕竟我赢先。\par

\section{临江仙}
\textbf{蔡伸}\par\vspace{0.5em}
仙品不同桃李艳，移来月窟云乡。\par
幽姿绰约道家妆。\par
绿云堆髻，娇额半涂黄。\par
可但乍凉风月下，饶伊独占秋光。\par
雨中别有恼人香。\par
错教萧史，肠断忆巫阳。\par

\section{临江仙}
\textbf{蔡伸}\par\vspace{0.5em}
琪树鸾栖花露重，依稀兰洞风光。\par
玉人相对自生凉。\par
翠鬟琼佩，绰约蕊珠妆。\par
宝瑟声沈清梦觉，夜阑明月幽窗。\par
可堪襟袂惹馀香。\par
断云残雨，何处认高唐。\par

\section{临江仙}
\textbf{蔡伸}\par\vspace{0.5em}
记得南楼三五夜，曾听凤管昭华。\par
曾前此际重兴嗟。\par
素娥端有恨，烟霭等闲遮。\par
珍重主人留客意，厌厌缓引流霞。\par
夜间银汉淡天涯。\par
亭亭丹桂现，耿耿玉绳斜。\par

\section{临江仙}
\textbf{蔡伸}\par\vspace{0.5em}
青润奇峰名韫玉，温其质并琼瑶。\par
中分瀑布泻云涛。\par
双峦呈翠色，气象两相高。\par
珍重幽人诚好事，绿窗聊助风骚。\par
寄言俗客莫相嘲。\par
物轻人意重，千里赠鹅毛。\par

\section{鹧鸪天}
\textbf{蔡伸}\par\vspace{0.5em}
脉脉柔情不自持。\par
浅颦轻笑百般宜。\par
尊前唱歇黄金缕，一点春愁入翠眉。\par
流蕙盼，捧瑶卮。\par
借君歌扇写新诗。\par
浮花谩说惊郎目，不似东风第一枝。\par

\section{惜奴娇}
\textbf{蔡伸}\par\vspace{0.5em}
隔阔多时，算彼此、难存济。\par
咫尺地、千山万水。\par
眼眼相看，要说话、都无计。\par
只是。\par
唱曲儿、词中认意。\par
雪意垂垂，更刮地、寒风起。\par
怎禁这几夜意。\par
未散痴心，便指望、长偎倚。\par
只替。\par
那火桶儿、与奴暖被。\par

\section{行香子}
\textbf{蔡伸}\par\vspace{0.5em}
珠露初零。\par
天宇澄明。\par
正闲阶、皎月亭亭。\par
更阑人静，烟敛风清。\par
更井边桐，一叶叶，做秋声。\par
斗帐鸾屏。\par
翠被华茵。\par
梦回时、酒力初醒。\par
绿云堆枕，红玉生春。\par
且打叠起，龙牙簟，竹夫人。\par

\section{一翦梅}
\textbf{蔡伸}\par\vspace{0.5em}
堆枕乌云堕翠翘。\par
午梦惊回，满眼春娇。\par
嬛嬛一袅楚宫腰。\par
那更春来，玉减香消。\par
柳下朱门傍小桥。\par
几度红窗，误认鸣镳。\par
断肠风月可怜宵。\par
忍使恹恹，两处无聊。\par

\section{一翦梅}
\textbf{蔡伸}\par\vspace{0.5em}
高宴华堂夜向阑。\par
急管飞霜，羯鼓声乾。\par
仙人掌上水晶盘。\par
回按凌波，舞袖弓弯。\par
曲罢凝娇整翠鬟。\par
玉笋持杯，巧笑嫣然。\par
为君一醉倒金船。\par
只恐醒来，人隔云山。\par

\section{一翦梅}
\textbf{蔡伸}\par\vspace{0.5em}
夜永虚堂烛影寒。\par
斗转春来，又是明年。\par
异乡怀抱只凄然。\par
尊酒相逢且自宽。\par
天际孤云云外山。\par
梦绕觚稜，日下长安。\par
功名已觉负初心，羞对菱花，绿鬓成斑。\par

\section{六么令}
\textbf{蔡伸}\par\vspace{0.5em}
梅英飘雪，弱柳弄新绿。\par
泠泠画桥流水，风静波如縠。\par
长记扁舟共载，偶近旗亭宿。\par
渺云横玉。\par
鸳鸯枕上，听彻新翻数般曲。\par
此际魂清梦冷，绣被香芬馥。\par
因念多感情怀，触处伤心目。\par
自是今宵独寐，怎不添愁蹙。\par
如今心足。\par
风前月下，赖有斯人慰幽独。\par

\section{镇西・小镇西犯}
\textbf{蔡伸}\par\vspace{0.5em}
秋风吹暗雨，重衾寒透。\par
伤心听、晓钟残漏。\par
凝情久。\par
记红窗夜雪，促膝围炉，交杯劝酒。\par
如今顿孤欢偶。\par
念别后。\par
菱花清镜里。\par
眉峰暗斗。\par
想标容、怎禁销瘦。\par
忍回首。\par
但云笺妙墨，鸳锦啼妆，依然似旧。\par
临风泪沾襟袖。\par

\section{看花回}
\textbf{蔡伸}\par\vspace{0.5em}
夜久凉生，庭院漏声频促。\par
念昔胜游旧地，对画阁层峦，雨馀烟簇。\par
新诗暗藏小字，霜刀刊翠竹。\par
携素手、细绕回塘，芰荷香里彩鸳宿。\par
别后想、香销腻玉。\par
带围减、钏宽金粟。\par
虽有鳞鸿锦素，柰事与心违，佳期难卜。\par
拟解愁肠万结，唯凭尊酒绿。\par
望天涯断魂处，醉拍阑干曲。\par

\section{诉衷情}
\textbf{蔡伸}\par\vspace{0.5em}
亭亭秋水玉芙蓉。\par
天际水浮空。\par
碧云望中空暮，人在广寒宫。\par
双缕枕，曲屏风。\par
小房栊。\par
可怜今夜，明月清风，无计君同。\par

\section{浪淘沙}
\textbf{蔡伸}\par\vspace{0.5em}
楼下水潺湲。\par
楼外屏山。\par
淡烟笼月晚凉天。\par
曾共玉人携素手，同倚阑干。\par
云散梦难圆。\par
幽恨绵绵。\par
旧游重到忍重看。\par
负你一生多少泪，月下花前。\par

\section{如梦令}
\textbf{蔡伸}\par\vspace{0.5em}
人静重门深亚。\par
朱阁画帘高挂。\par
人与月俱圆，月色波光相射。\par
潇洒。\par
潇洒。\par
人月长长今夜。\par

\section{如梦令}
\textbf{蔡伸}\par\vspace{0.5em}
今夜行云何处。\par
还是月华当午。\par
倚遍曲阑桥，望断锦屏归路。\par
空去。\par
空去。\par
梦到绿窗朱户。\par

\section{愁倚阑・春光好}
\textbf{蔡伸}\par\vspace{0.5em}
伤春晚，送春归。\par
步云溪。\par
绿叶同心双小字，记曾题。\par
楼外红日平西。\par
长亭路、烟草萋萋。\par
云雨不成新梦后，倚阑时。\par

\section{愁倚阑・春光好}
\textbf{蔡伸}\par\vspace{0.5em}
天如水，月如钩。\par
正新秋。\par
月影参差人窈窕，小红楼。\par
如今往事悠悠。\par
楼前水、肠断东流。\par
旧物忍看金约腕，玉搔头。\par

\section{愁倚阑・春光好}
\textbf{蔡伸}\par\vspace{0.5em}
一番雨，一番凉。\par
夜初长。\par
满院蛩吟人不寝，月侵廊。\par
木犀微绽幽芳。\par
西风透、窈窕红窗。\par
恰似个人鸳被里，玉肌香。\par

\end{document}